\documentclass[11pt]{article}
\usepackage[margin=1in]{geometry}
\usepackage[T1]{fontenc}
\usepackage[utf8]{inputenc}
\usepackage{lmodern}
\usepackage{microtype}
\usepackage{amsmath}
\usepackage{booktabs}
\usepackage{enumitem}
\usepackage{hyperref}

\title{GPQA Prompt-Profile Predictor Recommendation After the First \texttt{s\_0.9} Pilot}
\author{}
\date{2026-03-21}

\begin{document}
\maketitle

\paragraph{Bottom line.}
The next predictor should still be a prompt-level terminal-statistics predictor
from prefill activations, not another binary loop classifier. But after the
first real 2026-03-21 GPQA pilot, I would \emph{not} keep
\texttt{s\_0.9 = P(L / E >= 0.9)} as the sole first head. On this actual
\texttt{GPQA}, \texttt{temperature=0.2}, \texttt{n=10}, \texttt{max\_tokens=30000}
surface, \texttt{s\_0.9} collapsed exactly to \texttt{p(cap-hit)} prompt by
prompt, so it did not deliver the intended cap-robust tail target. The next
useful head should therefore be a denser terminal statistic:
\begin{itemize}[leftmargin=1.5em]
\item \textbf{ship next:} \texttt{mean\_relative\_length = E[L / E]}, and
\item \textbf{optional event-style companion:} a lower-threshold tail event such
as \texttt{s\_0.75} rather than \texttt{s\_0.9}.
\end{itemize}
Keep \texttt{p(max\_length\_hit)}, \texttt{p(loop)}, and correctness as
diagnostic outputs, not the headline loss.

\paragraph{Why the original \texttt{s\_0.9} idea was reasonable.}
The 2026-03-16 common-policy rollout-stat bundle on GPQA under
\texttt{Qwen/Qwen3-1.7B}, \texttt{temperature=0.2}, and \texttt{num\_generations=10}
shows that the pathology lives in the long right tail rather than in a single
binary cap-hit event:
\begin{itemize}[leftmargin=1.5em]
\item rollout success: \texttt{683/1980 = 34.5\%},
\item looped rollouts: \texttt{325/1980 = 16.4\%},
\item max-length hits: \texttt{137/1980 = 6.9\%},
\item cap-hit conditional on loop: \texttt{137/325 = 42.2\%},
\item loop conditional on cap-hit: \texttt{137/137 = 100.0\%},
\item average generation length: \texttt{9.7k} tokens,
\item average looped generation length: \texttt{23.5k} tokens.
\end{itemize}
So cap-hit is a precise fixed-policy alarm, but it only captures about
\texttt{42\%} of looped GPQA rollouts. That is why a tail-profile target was the
right initial direction: raw cap-hit alone discards much of the long-tail
behavior and is too tied to the chosen ceiling.

\paragraph{What the first real pilot changed.}
The resumed 2026-03-21 GPQA pilot completed end to end on a \texttt{32 / 16}
prompt split and produced the first actual \texttt{s\_0.9} probe artifact.
The important findings are:
\begin{itemize}[leftmargin=1.5em]
\item \texttt{s\_0.9} equals \texttt{p(cap-hit)} for \emph{every} prompt in this
pilot's archive; there were no prompts with substantial mass in the
\texttt{0.9E} to \texttt{E} band that avoided exact cap hits.
\item The target is very sparse: train target mean \texttt{0.0531}, test target
mean \texttt{0.1063}, with only \texttt{1/32} train prompts and \texttt{2/16}
test prompts at or above \texttt{0.5}.
\item The trained ensemble does not yet beat a constant baseline:
\begin{itemize}[leftmargin=1.5em]
\item probe eval Brier \texttt{0.03456} vs constant-baseline Brier
\texttt{0.03434},
\item probe eval MAE \texttt{0.13843} vs constant-baseline MAE
\texttt{0.13516},
\item eval Spearman \texttt{-0.2799}.
\end{itemize}
\item Prompt length still carries moderate signal into the target
(\texttt{Spearman(prompt\_len, s\_0.9) = 0.43} on test), so the leakage check
remains necessary.
\end{itemize}
This is enough to say that, on the current GPQA/\texttt{30k} surface,
\texttt{s\_0.9} is still too close to cap-hit to be the best first objective.

\paragraph{Concrete build.}
The first useful predictor should be:
\begin{enumerate}[leftmargin=1.5em]
\item \textbf{Dataset.} Same-source, prompt-disjoint GPQA split with the repaired
multiple-choice prompt formatter. Keep the decode policy fixed at
\texttt{temperature=0.2}, \texttt{num\_generations=10},
\texttt{max\_model\_len=40960}, and the current long rollout cap.
\item \textbf{Inputs.} Prefill activations only, using the last prompt token.
Do not concatenate all layers into one input vector. Instead compare:
\begin{itemize}[leftmargin=1.5em]
\item one last-layer MLP probe, and
\item one MLP per layer with probability aggregation by \texttt{mean\_prob}.
\end{itemize}
\item \textbf{Loss.} Use a dense terminal-statistics head first:
\begin{itemize}[leftmargin=1.5em]
\item regression on \texttt{mean\_relative\_length} for the next run, or
\item if an event target is preferred, lower the tail threshold toward the actual
GPQA loop-length regime (roughly \texttt{23.5k / 30k \(\approx\) 0.78}), for
example \texttt{s\_0.75} rather than \texttt{s\_0.9}.
\end{itemize}
Keep correctness out of the first loss so the model does not mostly learn task
difficulty.
\item \textbf{Evaluation.} Use probability-target metrics:
\texttt{Brier}, \texttt{MAE}, \texttt{Spearman}, and top-bucket capture for
probability targets, or \texttt{MSE}/\texttt{MAE}/\texttt{Spearman} for the
regression head. Report \texttt{p(max\_length\_hit)}, \texttt{p(loop)}, and
correctness only as downstream diagnostics on the same prompts.
\item \textbf{Baselines.} Hard-benchmark against prompt token count and effective
budget $E$ alone. The probe is only interesting if it beats those leakage-style
baselines.
\end{enumerate}

\paragraph{Recommended next sequence.}
\begin{enumerate}[leftmargin=1.5em]
\item Reuse the finished GPQA rollout archive to train a \texttt{mean\_relative\_length}
head on the same prompts and same prefill features.
\item If an event-style target is still wanted, rebuild only the labels from the
same archive at a lower threshold such as \texttt{s\_0.75}.
\item Keep \texttt{p(max\_length\_hit)} and \texttt{p(loop)} as readout
diagnostics rather than the main objective.
\end{enumerate}

\paragraph{Same-archive regression check.}
That first follow-up is already done on the exact same \texttt{32 / 16} prompts,
with no second rollout. The denser target is directionally better than
\texttt{s\_0.9} but still not good enough to call the predictor useful:
\begin{itemize}[leftmargin=1.5em]
\item best \texttt{mean\_relative\_length} run: eval MSE \texttt{0.0503}, eval
MAE \texttt{0.1549}, eval Spearman \texttt{+0.0647};
\item constant baseline on the same test targets: MSE \texttt{0.0311}, MAE
\texttt{0.1342};
\item prompt-length-only linear fit on the same split: MSE \texttt{0.0422}, MAE
\texttt{0.1425}.
\end{itemize}
So the target swap helps the rank correlation sign, but it does not yet beat
very weak baselines on a \texttt{48}-prompt pilot. The practical implication is
that the next round should keep the denser target \emph{and} grow prompt count,
instead of treating this micro-pilot as evidence that the predictor is already
working.

\paragraph{Current pilot status.}
The resumed 2026-03-21 GPQA pilot on GPUs \texttt{0-2} finished cleanly in
\texttt{37m22s}. It produced the full dataset archive, diagnostics, and probe
checkpoints, so the blocker is no longer runtime feasibility. The blocker is now
pilot scale plus objective choice: the fixed architecture and in-distribution
setup are workable, but \texttt{s\_0.9} at this cap is not yet a useful target,
and the denser same-archive regression head still needs more prompts before it
can be taken seriously.

\end{document}
